% GENERATED FILE. Edit canonical lessons, then run npm run generate:books.
% canonical-source-hash: fnv1a64:ad6480712c198ee2
% canonical-lessons: ZH-C367-wuhui, ZH-C367-liaojie, ZH-W367-jie-untie, ZH-C367-lijie, ZH-C367-jianjie, ZH-C367-kexiao

\chapter{I See What You Mean}
\label{ch:zh-i-see-what-you-mean}
\hlchaptermodality{367}

\begin{chapteropening}
\textbf{By the end of this chapter:} \emph{I can say I understand someone, know them well, and find an idea ridiculous.}

Complete the last lesson of chapter 367: I can say I understand someone, know them well, and find an idea ridiculous.
\end{chapteropening}

\section[wùhuì]{wùhuì — to get the wrong idea; to misunderstand}

\label{lesson:ZH-C367-wuhui}



\begin{quote}
Before the new one: say the Mandarin for to give up, then the Mandarin for to be worth it.
\end{quote}

\subsection*{You'll want to know: wùhuì}
\emph{wùhuì} — \textquotedblleft{}to get the wrong idea; to misunderstand\textquotedblright{}.

\begin{quote}
\emph{nǐ wùhuì wǒ de yìsi le} — You've misunderstood me.
\end{quote}

Its characters come later; for now it is a word for the ear.

\begin{cousinweb}
\emph{wùhuì} is built from \emph{wù} and \zh{会} 'to meet', both inherited from Middle Chinese, \emph{nguH} and \emph{hwajH} (Guangyun, in Baxter's transcription), which glosses the first 'an error' and the second 'to come together'.
\end{cousinweb}

\begin{grammarlens}[title={What you've built}]
One more word for giving reasons.
\end{grammarlens}

\subsection*{Your turn}
\begin{itemize}
\raggedright
  \item \emph{Say it:} \emph{wùhuì}
  \item \emph{Say it:} \emph{wùhuì}, once more
  \item \emph{Say it:} \emph{nǐ wùhuì wǒ de yìsi le} — You've misunderstood me.
  \item \emph{Recall:} say the Mandarin for to give up, then the Mandarin for to be worth it, then say \emph{wùhuì} again
\end{itemize}

\begin{tcolorbox}[breakable,colback=teal!4,colframe=teal!35!black,title={Before you move on}]
What does wùhuì mean? (\textquotedblleft{}To get the wrong idea; to misunderstand\textquotedblright{}.) Say it once more.
\end{tcolorbox}
\section[\texorpdfstring{\zh{了解} (liǎojiě)}{了解 (liǎojiě)}]{\zh{了解} (liǎojiě) — to know in depth; to find out about}

\label{lesson:ZH-C367-liaojie}



\begin{quote}
Before the new one: say the Mandarin for to be worth it, then the Mandarin for to get the wrong idea.
\end{quote}

\subsection*{You'll want to know: \zh{了解}}
\textbf{\zh{了解}} — \emph{liǎojiě} — \textquotedblleft{}to know in depth; to find out about\textquotedblright{}.

\begin{quote}
\textbf{\zh{我很了解他}} \emph{wǒ hěn liǎojiě tā} — I know him well.
\end{quote}

The new character, in pieces: \zh{解} — thirteen strokes in three pieces: on the left a horn piece, seven strokes; on the top right a knife piece, two strokes; on the bottom right a cow piece, four strokes. The next lesson puts it in your hand.

Two third tones meet in this word: liǎo is said with a rising tone before jiě, and the spelling keeps both marks.

\begin{cousinweb}
\textbf{\zh{了解}} (\emph{liǎojiě}) is built from \zh{了} and \zh{解}, both inherited from Middle Chinese, \emph{lewX} and \emph{keaX} (Guangyun, in Baxter's transcription), which glosses the first 'clever' and 'finished', and the second 'to explain' and 'to loosen'.
\end{cousinweb}

\begin{grammarlens}[title={What you've built}]
One more word for giving reasons.
\end{grammarlens}

\subsection*{Your turn}
\begin{itemize}
\raggedright
  \item \emph{Say it:} \emph{liǎojiě}
  \item \emph{Say it:} \emph{liǎojiě}, once more
  \item \emph{Say it:} \emph{wǒ hěn liǎojiě tā} — I know him well.
  \item \emph{Recall:} say the Mandarin for to be worth it, then the Mandarin for to get the wrong idea, then say \emph{liǎojiě} again
\end{itemize}

\begin{tcolorbox}[breakable,colback=teal!4,colframe=teal!35!black,title={Before you move on}]
What does \zh{了解} mean? (\textquotedblleft{}To know in depth; to find out about\textquotedblright{}.) Say it once more.
\end{tcolorbox}
\section[\texorpdfstring{\zh{解} (jiě)}{解 (jiě)}]{\zh{解} — a horn, a knife and a cow}

\label{lesson:ZH-W367-jie-untie}



\begin{quote}
The word first, then the character inside it.

\begin{itemize}
\raggedright
  \item \emph{Recall:} say \emph{liǎojiě} — \textbf{R1}, one lesson back, before it decays
\end{itemize}

\begin{quote}
\zh{了} \zh{见}
\end{quote}

\emph{Read it:} each one aloud, with nothing said first — \emph{le}, \emph{jiàn}
\end{quote}

\subsection*{Script — a horn, a knife and a cow}
\hlblockfigure{\detokenize{figures/ZH-W367-jie-untie-filmstrip.pdf}}{How \zh{解} is written, stroke by stroke}

\begin{quote}
\zh{解}
\end{quote}

\textbf{Three pieces: the horn, then the knife, then the cow.} Count three, not thirteen.

\textbf{The horn piece, seven strokes}, on the left: a short left-falling stroke, a horizontal that turns down to the left, a long left-falling stroke down the left side, a horizontal that turns down and hooks, two short horizontals inside, then a vertical down the middle.

\textbf{The knife piece, two strokes}, on the top right: a horizontal that turns down and hooks, then a left-falling stroke inside it.

\textbf{The cow piece, four strokes}, on the bottom right: a short left-falling stroke, a short horizontal, a long horizontal, then a vertical down through them.

7 + 2 + 4 = thirteen.

\subsection*{Your turn}
\begin{itemize}
\raggedright
  \item \emph{Write it:} \zh{解}, all thirteen, with the model in front of you
  \item \emph{Cover:} write it again from memory, and count the pieces, not the strokes
\end{itemize}

\begin{tcolorbox}[breakable,colback=teal!4,colframe=teal!35!black,title={Before you move on}]
How many pieces in \zh{解}? (\textbf{Three}: the horn, the knife, the cow.) Which piece comes first? (\textbf{The horn, on the left}.) How many strokes does the knife have? (\textbf{Two}.)

Source: \href{https://raw.githubusercontent.com/chanind/hanzi-writer-data/68d10a4b21150cae5e1ebbd223eed289cf32d90c/data/\%E8\%A7\%A3.json}{Hanzi Writer Data \zh{解}, pinned stroke paths and medians}.
\end{tcolorbox}
\section[\texorpdfstring{\zh{理解} (lǐjiě)}{理解 (lǐjiě)}]{\zh{理解} (lǐjiě) — to understand someone's reasons; to sympathize}

\label{lesson:ZH-C367-lijie}



\begin{quote}
Before the new one: say the Mandarin for to get the wrong idea, then the Mandarin for to know in depth.
\end{quote}

\subsection*{You'll want to know: \zh{理解}}
\textbf{\zh{理解}} — \emph{lǐjiě} — \textquotedblleft{}to understand someone's reasons; to sympathize\textquotedblright{}.

\begin{quote}
\textbf{\zh{因为我也当过学生所以我很理解你}} \emph{yīnwèi wǒ yě dāng guo xuésheng suǒyǐ wǒ hěn lǐjiě nǐ} — I was a student once too, so I really do understand.
\end{quote}

Two third tones meet in this word: lǐ is said with a rising tone before jiě, and the spelling keeps both marks.

\begin{cousinweb}
\textbf{\zh{理解}} (\emph{lǐjiě}) is built from \zh{理} 'order' and \zh{解} 'to untie', both inherited from Middle Chinese, \emph{liX} and \emph{keaX} (Guangyun, in Baxter's transcription), which glosses the first 'to put in order' and the second 'to explain'.
\end{cousinweb}

\begin{grammarlens}[title={What you've built}]
One more word for giving reasons.
\end{grammarlens}

\subsection*{Your turn}
\begin{itemize}
\raggedright
  \item \emph{Say it:} \emph{lǐjiě}
  \item \emph{Say it:} \emph{lǐjiě}, once more
  \item \emph{Say it:} \emph{yīnwèi wǒ yě dāng guo xuésheng suǒyǐ wǒ hěn lǐjiě nǐ} — I was a student once too, so I really do understand.
  \item \emph{Recall:} say the Mandarin for to get the wrong idea, then the Mandarin for to know in depth, then say \emph{lǐjiě} again
\end{itemize}

\begin{tcolorbox}[breakable,colback=teal!4,colframe=teal!35!black,title={Before you move on}]
What does \zh{理解} mean? (\textquotedblleft{}To understand someone's reasons; to sympathize\textquotedblright{}.) Say it once more.
\end{tcolorbox}
\section[\texorpdfstring{\zh{见解} (jiànjiě)}{见解 (jiànjiě)}]{\zh{见解} (jiànjiě) — insight; considered view}

\label{lesson:ZH-C367-jianjie}



\begin{quote}
Before the new one: say the Mandarin for to know in depth, then the Mandarin for to understand someone's reasons.
\end{quote}

\subsection*{You'll want to know: \zh{见解}}
\textbf{\zh{见解}} — \emph{jiànjiě} — \textquotedblleft{}insight; considered view\textquotedblright{}.

\begin{quote}
\textbf{\zh{他的见解很有意思}} \emph{tā de jiànjiě hěn yǒu yìsi} — His views are interesting.
\end{quote}

\begin{cousinweb}
\textbf{\zh{见解}} (\emph{jiànjiě}) is built from \zh{见} 'to see' and \zh{解} 'to untie', both inherited from Middle Chinese, \emph{kenH} and \emph{keaX} (Guangyun, in Baxter's transcription), which glosses the first 'to see' and the second 'to explain'.
\end{cousinweb}

\begin{grammarlens}[title={What you've built}]
One more word for giving reasons.
\end{grammarlens}

\subsection*{Your turn}
\begin{itemize}
\raggedright
  \item \emph{Say it:} \emph{jiànjiě}
  \item \emph{Say it:} \emph{jiànjiě}, once more
  \item \emph{Say it:} \emph{tā de jiànjiě hěn yǒu yìsi} — His views are interesting.
  \item \emph{Recall:} say the Mandarin for to know in depth, then the Mandarin for to understand someone's reasons, then say \emph{jiànjiě} again
\end{itemize}

\begin{tcolorbox}[breakable,colback=teal!4,colframe=teal!35!black,title={Before you move on}]
What does \zh{见解} mean? (\textquotedblleft{}Insight; considered view\textquotedblright{}.) Say it once more.
\end{tcolorbox}
\section[\texorpdfstring{\zh{可笑} (kěxiào)}{可笑 (kěxiào)}]{\zh{可笑} (kěxiào) — ridiculous; laughable}

\label{lesson:ZH-C367-kexiao}



\begin{quote}
Before the new one: say the Mandarin for to understand someone's reasons, then the Mandarin for an insight.
\end{quote}

\subsection*{You'll want to know: \zh{可笑}}
\textbf{\zh{可笑}} — \emph{kěxiào} — \textquotedblleft{}ridiculous; laughable\textquotedblright{}.

\begin{quote}
\textbf{\zh{我觉得这个问题很可笑}} \emph{wǒ juéde zhège wèntí hěn kěxiào} — I think that question is ridiculous.
\end{quote}

\begin{cousinweb}
\textbf{\zh{可笑}} (\emph{kěxiào}) is built from \zh{可} 'can' and \zh{笑} 'to laugh', both inherited from Middle Chinese, \emph{khaX} and \emph{sjewH} (Guangyun, in Baxter's transcription), which glosses the first 'to permit' and the second 'glad'.
\end{cousinweb}

\begin{grammarlens}[title={What you've built}]
One more word for giving reasons.
\end{grammarlens}

\subsection*{Your turn}
\begin{itemize}
\raggedright
  \item \emph{Say it:} \emph{kěxiào}
  \item \emph{Say it:} \emph{kěxiào}, once more
  \item \emph{Say it:} \emph{wǒ juéde zhège wèntí hěn kěxiào} — I think that question is ridiculous.
  \item \emph{Recall:} say the Mandarin for to understand someone's reasons, then the Mandarin for an insight, then say \emph{kěxiào} again
\end{itemize}

\begin{tcolorbox}[breakable,colback=teal!4,colframe=teal!35!black,title={Before you move on}]
What does \zh{可笑} mean? (\textquotedblleft{}Ridiculous; laughable\textquotedblright{}.) Say it once more.
\end{tcolorbox}
